\documentclass{ximera}

\title{Power Series Activity}
\author{MATH 426: Calculus II}

\begin{document}
\begin{abstract}
   Working with the peers in your group, solve the following problems. Make sure to show and justify all your work. Make sure everyone in the group understands the solution and participates. Be prepared to report your answers to the whole class. 
\end{abstract}
\maketitle

\textbf{Purpose:} In this activity you will practice finding the interval of convergence of power series. Rather than work through disconnected examples, this activity will use our existing knowledge of geometric series to create a set of example problems which smoothly increase in technical difficulty so that we can practice the basic task of finding the interval of convergence.

\textbf{Pre-requisites:} Familiarity with geometric series, $p$-series and the ratio test. A brief introduction to the alternating series test will be helpful on the later questions.\\
\end{tcolorbox}


\begin{exercise}
To generate our first example power series, we return to the result (proven in lecture) about the value of geometric series.
    \begin{align}
	    \sum_{n=0}^\infty (r)^n = \frac{1}{1-r}, \hspace{.2in} \text{ when }|r|<1\nonumber\end{align}
Without proof we will try to interpret this as an equation between functions:
    \begin{align}
	    \sum_{n=0}^\infty x^n = \frac{1}{1-x} \nonumber
    \end{align}
The left hand side of this equation fits our definition of a power series, with \textbf{all} of the coefficients equal to 1.\\  

From here we will attempt to find the interval of convergence of the power series.\\

To find the interval of convergence we employ the following sequence of steps:
\begin{enumerate}
    \item Apply the ratio test to the terms of our series. (While performing this convergence test, we treat $x$ as a constant parameter at first.)

    \item When we get to the interpretation stage of the ratio test (i.e. after we finish evaluating the limit), we check to see if we can change the outcome of the test based on the value of $x$. Note: If there is a value of $x$ where the limit becomes greater than 1, the radius of convergence is finite.
    
    \item If we find a finite radius of convergence, then we substitute the inconclusive values of $x$ into the series and use {\bf{other}} convergence tests to determine whether the series converges at the endpoints. \\
    Note: Because the inconclusive values of $x$ will result in one alternating series and one non-alternating series, it is most efficient to analyze endpoint convergence by checking the non-alternating series first.  If this series converges, it means the alternating series absolutely converges (and thus converges itself).  
\end{enumerate}
\end{exercise}


For problems 2-5, we will continue to find intervals of convergence for power series.  Try one example from each problem to start, then return to the others as you have time.

\begin{exercise}
Repeat the process to find the interval of convergence of more complicated series built from purely geometric parts. 
\begin{align}
	\sum_{n=0}^\infty, \left(\frac{x}{3}\right)^n, \hspace{.5in} \sum_{n=0}^\infty \left(-x\right)^n, \hspace{.5in} \sum_{n=0}^\infty \left(4x\right)^n,\nonumber
\end{align}
\textit{Hint: All three of these examples have open intervals of convergence: $(a,b)$. You should find these series diverges at both endpoints.}
\end{exercise}

\begin{exercise}
Repeat the process and try to find the interval of convergence of more complicated series with a harmonic feature thrown in. 
\begin{align}
	\sum_{n=1}^\infty, \frac{1}{n}\left(\frac{x}{2}\right)^n, \hspace{.5in} \sum_{n=1}^\infty \frac{\left(-x\right)^n}{n}, \hspace{.5in} \sum_{n=1}^\infty \frac{\left(5x\right)^n}{n},\nonumber
\end{align}
\textit{Hint: All three of these examples have `half-open' intervals of convergence: $[a,b)$ or $(a,b]$. You should find the power series converges at exactly one endpoint.}

\end{exercise}

\begin{exercise}
Repeat the process and try to find the interval of convergence of more complicated series with a convergent $p$-series term thrown in.  
\begin{align}
	\sum_{n=1}^\infty, \frac{x^n}{n^2}, \hspace{.5in} \sum_{n=1}^\infty \frac{\left(-x\right)^n}{n^3}, \hspace{.5in} \sum_{n=0}^\infty \frac{\left(5x\right)^n}{(n+1)^2},\nonumber
\end{align}
\textit{Hint: All three of these examples have closed intervals of convergence: $[a,b]$ or $[a,b]$. You should find the power series converges at both endpoints.}
\end{exercise}

\begin{exercise}
(Time permitting:) Graph the function $\frac{1}{1-x}$ and note any links you see between the function and the radius/interval of convergence.
\end{exercise}

\begin{exercise}
(Time permitting): As a group see if you can create an examples of {\bf{shifted}} power series:
\begin{align}
	\sum_{n=1}^\infty f(n)(x-c)^n\nonumber
\end{align}
with enough flexibility to `customize' the interval of convergence. $(a,b)$, $[a,b)$, $(a,b]$, or $[a,b]$.

\end{exercise}


\end{document}



